\section{Introduction}
\label{sec:introduction}

Large language models (LLMs) answer factual questions fluently, but fluency is a poor guide to correctness \cite{brown2020gpt3,ji2023hallucination}. A model may repeat an outdated fact, confuse two entities with similar names, or state something plausible that no source supports, and the reader cannot tell which part of the answer rests on evidence. Two families of methods attack this weakness from different sides. Retrieval-augmented generation places passages retrieved from an external corpus in front of the model, so that the answer can be grounded in text outside the model \cite{lewis2020rag,guu2020realm,izacard2021fid}. Multi-agent debate gives the same question to several instances of a model and lets them read and criticise each other's answers over a few rounds, on the premise that one instance catches the errors of another \cite{du2024debate,liang2024divergent,chan2024chateval}.

MADKE \cite{wang2025madke} joins the two ideas. Each question comes with a shared pool of Wikipedia and Google passages. In every round, each debating agent selects at most three passages from the pool, answers with them in view and reads the answers of the other agents. A judge agent decides after the round whether the agents agree, and a summariser agent writes the final answer once they agree or the round budget is used up. The method appeared in \emph{Neurocomputing} and is reported to outperform single-agent prompting, chain-of-thought prompting \cite{wei2022cot}, self-consistency \cite{wang2023sc}, retrieval-augmented generation and plain debate on TriviaQA, Natural Questions, HotpotQA, 2WikiMultihopQA and FEVER with three backbone models.

Other researchers can use a method of this kind only through its implementation. Before MADKE was extended in this study, its public code \cite{madke2025code} was read line by line and its logic was run again under controlled replies. The code does not execute the method as the paper describes it. The call of the summariser passes the debate transcript and the question in the wrong order, so that in every final prompt the line marked ``Question:'' holds the replies of the agents. The parser for the passages an agent selects accepts only brackets that hold a single number: a selection written as one list, such as \texttt{[2, 0, 7]}, delivers no passage at all, and nothing in the output says so. The judge counts as agreeing only if it writes the exact string \texttt{[Yes]}; with \texttt{[yes]} or \texttt{["Yes"]} the debate goes on. The script that computes the consistency statistics never counts an agreement reached in the final round. A failed model call hands back its error message as if it were a reply, and the message is scored as an answer. The default pool is larger than the paper specifies, it can hold the same passage twice, and it is always listed with Wikipedia first. The output keeps the raw prompts and replies, but no structured record of which passages reached which agent and which of them support the final answer. Correctness, cost and auditability are all affected.

This study responds with MADKE-X. The debate itself is left untouched: the same agents, the same debate prompts, the same round budget. A difference between the two systems can therefore be attributed to what was added around the debate and not to a new debate design. Four components are added. A failure-aware interface layer makes every model call report success or failure, binds the summariser prompt by keyword, reads the selection of an agent from the end of its reply in any bracket format, reads the last verdict of the judge in any spelling, and computes the debate statistics from stored verdicts. An evidence governor builds the specified pool of five plus five passages, removes duplicates and orders the rest by reciprocal rank fusion \cite{cormack2009rrf} of the retrieval rank and a BM25 rank \cite{robertson2009bm25} computed from the question alone. A decision trace records the pool, the passages delivered in every round, the verdicts and the passages the summariser cites with $\{E\#\}$ markers, together with a lexical check of whether the answer occurs in a cited passage. A bounded verification step calls the summariser once more when that check fails and adopts the second answer only if it passes.

The evaluation separates what can be shown without a language model from what cannot. The five benchmark files shipped with MADKE are used unaltered and are verified by SHA-256 digests before every run. The order of the evidence, the size of the prompts and the cost of the added steps are properties of the data and the code, and they are measured directly on the \VNQuestions{} questions. Whether a pipeline does what it is specified to do is tested with deterministic scripted agents, which drive the release and MADKE-X through replies of known content so that every difference in behaviour is caused by the code. End-to-end accuracy needs a live model. For that comparison both systems answer the same questions with the same open-weight backbone under a paired protocol; its status is reported in Section~\ref{sec:live}.

The contributions are the following.
\begin{enumerate}
\item An inspection of the public MADKE implementation that documents thirteen deviations from the published method or weaknesses of its evaluation, each tied to a statement of the code and, where the code can be exercised offline, reproduced by a test (Section~\ref{sec:baseline}).
\item MADKE-X, a refinement of MADKE with four components, each defined by an equation whose symbols and effect on the results are stated (Section~\ref{sec:method}), together with a complexity analysis (Section~\ref{sec:complexity}).
\item A comparison protocol that re-runs the release under the same harness as MADKE-X and adds paired tests, bootstrap intervals, multiplicity correction and a power analysis, so that a claimed difference can be tested (Section~\ref{sec:design}).
\item Measurements on the unaltered data: the evidence governor lists answer-bearing passages earlier on all four question-answering datasets, the evidence listing shrinks by \VTokRedMean\%, and every answer carries a trace that reproduces what the pipeline did (Section~\ref{sec:offline}).
\item A strict separation between demonstrated and pending claims. Accuracy is reported only from live runs of both systems on identical questions, and no accuracy value is estimated or extrapolated (Section~\ref{sec:live}).
\end{enumerate}

The rest of the paper is organised as follows. Section~\ref{sec:related_work} reviews related work. Section~\ref{sec:baseline} describes MADKE and the findings on its implementation. Section~\ref{sec:method} presents MADKE-X with its equations, and Section~\ref{sec:complexity} analyses its cost. Section~\ref{sec:design} sets out the experimental design. Sections~\ref{sec:offline} and~\ref{sec:live} report the offline results and the live comparison, Section~\ref{sec:discussion} discusses them, Section~\ref{sec:threats} lists the threats to validity and Section~\ref{sec:conclusion} concludes.
